% =============================================================================
% 11-parallel-pipeline-hazards.tex -- Week 11: Flynn's classification,
% parallel processing, instruction pipelines, and pipeline hazards
%
% Citation discipline:
%   * Hamacher Ch. 5 and Ch. 10 are the syllabus anchors. The current
%     Hamacher2002 index does not page-verify these chapters, so citations stay
%     chapter-level and are presented as standard course treatment.
%   * Stallings is the 11th-edition PDF in the repository despite the retained
%     Stallings2015 key. Flynn appears only briefly in Ch. 20 Sec. 20.1; the
%     indexed pipeline-hazard treatment is Ch. 16 Sec. 16.4 (not the syllabus's
%     old Ch. 17-18 numbering).
%
% Compile from Slides/CS401 with Windows MiKTeX/XeLaTeX:
%   TEXINPUTS="../../Preambles;" xelatex -interaction=nonstopmode 11-parallel-pipeline-hazards.tex
%   BIBINPUTS="../..;" bibtex 11-parallel-pipeline-hazards
%   TEXINPUTS="../../Preambles;" xelatex -interaction=nonstopmode 11-parallel-pipeline-hazards.tex
%   TEXINPUTS="../../Preambles;" xelatex -interaction=nonstopmode 11-parallel-pipeline-hazards.tex
% =============================================================================

\documentclass[aspectratio=169]{beamer}
\input{header}
\coursecode{CS401}

\title{Parallel Processing and Pipeline Hazards}
\subtitle{Week 11 --- Flynn's Taxonomy, Instruction-Level Parallelism, and Safe Throughput}
\author[Dr. Auqib Hamid Lone]{%
  \textbf{Dr. Auqib Hamid Lone}\\
  \small Assistant Professor in Computer Science \& Engineering
}
\institute[GCET Ganderbal]{%
  \small\textbf{PCC CS-401: Computer Organization and Architecture}\\[0.3em]
  Department of Computer Science \& Engineering\\
  Government College of Engineering \& Technology (GCET), Ganderbal
}
\date[\today]{\small\today}

\begin{document}

\frame{\titlepage}

% =============================================================================
% ACT 0 --- RECAP, MOTIVATION, AND ROADMAP
% =============================================================================
\section{Recap and Motivation}

\begin{frame}{Where Week 10 Left Us}
  \begin{block}{One bottleneck kept changing shape}
    \begin{itemize}
      \item Week 9 used locality to hide DRAM latency behind cache.
      \item Week 10 used pages, frames, and storage to give a process a larger,
        protected address space than physical memory alone could provide.
      \item Both designs ask the same question: \key{how can the machine keep
        useful work moving while one operation is slow?}
    \end{itemize}
  \end{block}
  \begin{exampleblock}{Today's turn}
    What if the processor itself has several useful operations ready at once?
    Parallel organizations and pipelines try to overlap those operations, but
    overlap is safe only when resources and dependencies permit it.
  \end{exampleblock}
\end{frame}

\begin{frame}{Motivation: One Instruction at a Time Wastes Hardware}
  \begin{exampleblock}{A five-step instruction path}
    Suppose each instruction needs five stages: fetch, decode, execute, memory,
    and writeback. If one instruction occupies the whole path before the next
    starts, four stage resources sit idle most of the time.
  \end{exampleblock}
  \begin{itemize}
    \item A \key{pipeline} overlaps different instructions in different stages.
    \item After the pipeline fills, one completed instruction can emerge each
      cycle even though each instruction still takes several stages of latency.
    \item The promise is \key{throughput}, not magic reduction of every
      instruction's individual latency.
  \end{itemize}
  \muted{This is the standard pipeline motivation in Hamacher Ch.~10
  \cite{Hamacher2002_computer_organization} and Stallings 11e, \S16.4,
  PDF pp.~667--685 \cite{Stallings2015_computer_organization}.}
\end{frame}

\begin{frame}{Roadmap: From Many Processors to One Safe Stream}
  \begin{enumerate}
    \item Classify parallel machines by instruction streams and data streams.
    \item Distinguish data-level, thread-level, and instruction-level
      parallelism.
    \item Build a five-stage instruction pipeline around one running stream.
    \item Diagnose structural, data, and control hazards.
    \item Resolve hazards with scheduling, duplication, forwarding, stalls,
      and prediction.
    \item Compare design choices: depth, clock period, bubbles, and complexity.
  \end{enumerate}
  \vspace{0.15cm}
  \begin{alertblock}{Running instruction stream}
    \texttt{LW R1,0(R2)}; \texttt{ADD R3,R1,R4}; \texttt{SUB R5,R3,R6};
    \texttt{BEQ R5,R0,L}; \texttt{OR R7,R8,R9}
  \end{alertblock}
\end{frame}

\transitionslide{First classify the kind of parallel work the machine exposes.}

% =============================================================================
% ACT 1 --- FLYNN AND PARALLELISM
% =============================================================================
\section{Flynn's Classification}

\begin{frame}{Flynn's Question: How Many Streams Are Moving?}
  \begin{block}{The organizing idea}
    Classify a computer by the number of concurrent \key{instruction streams}
    and \key{data streams} it processes. This is a taxonomy of organization,
    not a ranking from ``bad'' to ``good.''
  \end{block}
  \begin{itemize}
    \item \key{Instruction stream:} the sequence of operations being issued.
    \item \key{Data stream:} the operands or data elements those operations
      consume.
    \item A single physical chip can support more than one mode at different
      levels: for example, multiple cores may each run a SISD-like stream while
      a vector unit executes SIMD operations.
  \end{itemize}
  \muted{Flynn's taxonomy is standard course treatment from Hamacher Ch.~5
  \cite{Hamacher2002_computer_organization}; Stallings 11e's brief
  parallel-organization discussion is in \S20.1
  \cite{Stallings2015_computer_organization}.}
\end{frame}

\begin{frame}{The Four Flynn Classes}
  \scriptsize
  \begin{center}
    \begin{tabular}{@{}p{1.5cm}p{2.9cm}p{2.9cm}p{4.3cm}@{}}
      \toprule
      \textbf{Class} & \textbf{Instruction streams} & \textbf{Data streams} & \textbf{Typical interpretation} \\
      \midrule
      SISD & One & One & Conventional sequential core; one instruction acts on one data item \\
      SIMD & One & Many & One operation applied across a vector or image row \\
      MISD & Many & One & Rare and specialized; multiple operations inspect one stream \\
      MIMD & Many & Many & Multicore, multiprocessor, or cluster with independent work \\
      \bottomrule
    \end{tabular}
  \end{center}
  \vspace{0.15cm}
  \begin{alertblock}{Exam discipline}
    Count streams, then justify with the work pattern. ``Many ALUs'' alone does
    not establish SIMD; the instruction stream must be shared.
  \end{alertblock}
\end{frame}

\begin{frame}{Worked Classification: Three Workloads}
  \begin{exampleblock}{Classify before naming the hardware}
    \begin{enumerate}
      \item Add two 1024-element vectors using one vector instruction per chunk.
      \item Run four independent matrix tiles on four cores.
      \item Apply several independent filters to the same sensor stream.
    \end{enumerate}
  \end{exampleblock}
  \begin{itemize}
    \item (1) is \key{SIMD}: one operation pattern, many data elements.
    \item (2) is \key{MIMD}: independent instruction streams and data regions.
    \item (3) resembles \key{MISD} only under a strict stream definition; in
      practice it is often implemented as MIMD workers sharing an input buffer.
    \item The taxonomy describes the organization; the workload decides whether
      that organization is useful.
  \end{itemize}
\end{frame}

\begin{frame}{Parallelism Has Three Different Scales}
  \begin{block}{Do not collapse these into one word}
    \begin{itemize}
      \item \key{Data-level parallelism (DLP):} the same operation over many
        elements, as in vector addition or pixels.
      \item \key{Thread-level parallelism (TLP):} independent instruction
        sequences execute on separate cores or hardware threads.
      \item \key{Instruction-level parallelism (ILP):} one instruction stream
        contains independent instructions that the processor may overlap.
    \end{itemize}
  \end{block}
  \begin{itemize}
    \item Flynn's classes describe streams; DLP/TLP/ILP describe where useful
      concurrency comes from.
    \item The rest of this deck focuses on \key{ILP inside one instruction
      stream}, where hazards are the price of overlap.
  \end{itemize}
\end{frame}

\begin{frame}{Socratic Check: Which Parallelism Is Present?}
  \begin{alertblock}{Scenario}
    A processor has four cores. Each core runs a different loop iteration, and
    each loop iteration contains a five-stage pipeline. Which kinds of
    parallelism are present?
  \end{alertblock}
  \begin{itemize}
    \item \good{TLP:} four cores execute independent instruction streams.
    \item \good{ILP:} each core overlaps instructions within its own stream.
    \item \bad{Not automatically SIMD:} the cores may execute different
      instructions or different control paths.
    \item A good answer names the level and the evidence, rather than counting
      only the number of cores.
  \end{itemize}
\end{frame}

\transitionslide{Now overlap one instruction stream without violating its meaning.}

% =============================================================================
% ACT 2 --- PIPELINE MODEL
% =============================================================================
\section{Pipeline Processing}

\begin{frame}{The Five-Stage Pipeline Model}
  \begin{block}{A useful teaching model}
    \begin{tabular}{@{}lllll@{}}
      \textbf{IF} & \textbf{ID} & \textbf{EX} & \textbf{MEM} & \textbf{WB} \\
      Fetch & Decode/read & ALU/branch & Data memory & Register write
    \end{tabular}
  \end{block}
  \begin{itemize}
    \item Pipeline registers separate stages and hold intermediate results.
    \item Every stage advances on a clock edge; the slowest stage constrains
      the clock period.
    \item The same instruction visits all five stages, but different
      instructions occupy different stages simultaneously.
  \end{itemize}
  \muted{The stage names are a simplified course model; real processors may
  split, merge, repeat, or rename these functions.}
\end{frame}

\begin{frame}{Timing Example 1: Filling a Five-Stage Pipeline}
  \scriptsize
  \begin{exampleblock}{Four independent instructions}
    Assume no hazards and one cycle per stage. Compare sequential execution
    with a five-stage pipeline.
  \end{exampleblock}
  \begin{center}
    \begin{tabular}{@{}c|cccccccc@{}}
      \toprule
      & 1 & 2 & 3 & 4 & 5 & 6 & 7 & 8 \\
      \midrule
      $I_1$ & IF & ID & EX & MEM & WB & & & \\
      $I_2$ & & IF & ID & EX & MEM & WB & & \\
      $I_3$ & & & IF & ID & EX & MEM & WB & \\
      $I_4$ & & & & IF & ID & EX & MEM & WB \\
      \bottomrule
    \end{tabular}
  \end{center}
  \begin{itemize}
    \item Sequential: $4\times5=20$ stage-cycles. Pipelined: $5+(4-1)=8$.
    \item Speedup here is $20/8=2.5$, below the ideal five because the stream
      is finite and the pipeline must fill and drain.
  \end{itemize}
\end{frame}

\begin{frame}{Latency, Throughput, and the Ideal Limit}
  \begin{block}{Three quantities, three questions}
    \begin{itemize}
      \item \key{Latency:} how long from an instruction's IF to its WB?
      \item \key{Throughput:} how often can completed instructions emerge?
      \item \key{Speedup:} how does total completion time compare with a
        non-overlapped baseline?
    \end{itemize}
  \end{block}
  \begin{itemize}
    \item With $k$ balanced stages and $n$ independent instructions,
      $T_{pipe}\approx k+(n-1)$ cycles; sequential time is $nk$ cycles.
    \item For large $n$, ideal speedup approaches $k$, but only if stages are
      balanced and hazards, fills, drains, and overhead are negligible.
    \item A deeper pipeline can raise clock frequency while increasing branch
      penalty, register overhead, and recovery cost.
  \end{itemize}
\end{frame}

\begin{frame}{Pipeline Architecture: The Design Questions}
  \begin{enumerate}
    \item \key{Stage balance:} can work be divided so no stage dominates the
      clock period?
    \item \key{Pipeline registers:} what state must cross each boundary, and
      what setup/clock-to-q overhead is added?
    \item \key{Resource replication:} can instruction fetch and data access
      happen together, or must they share one memory port?
    \item \key{Dependence handling:} where can a value be forwarded, and when
      must the pipeline wait?
    \item \key{Control recovery:} how many wrong-path instructions are flushed
      after a branch is resolved?
  \end{enumerate}
  \muted{These are architecture choices, not merely implementation details:
  they determine CPI, area, power, and verification burden.}
\end{frame}

\begin{frame}{Socratic Check: What Does Pipelining Not Do?}
  \begin{alertblock}{Claim}
    ``A five-stage pipeline makes every instruction take one cycle.'' Is that
    statement correct?
  \end{alertblock}
  \begin{itemize}
    \item \bad{No.} An individual instruction still traverses five stages, so
      its latency is roughly five cycles in this model.
    \item \good{The valid claim:} after fill, independent instructions can
      complete at a rate of one per cycle.
    \item Stalls and flushes reduce the achieved throughput; they appear as
      bubbles or discarded work in the timing chart.
  \end{itemize}
\end{frame}

\transitionslide{The pipeline is fast only when its assumptions remain true.}

% =============================================================================
% ACT 3 --- HAZARDS
% =============================================================================
\section{Pipeline Hazards}

\begin{frame}{Hazard Taxonomy: Why Would Overlap Be Unsafe?}
  \begin{block}{A hazard is a condition that prevents the next instruction
  from executing in its intended cycle.}
    \begin{itemize}
      \item \key{Structural hazard:} two stages need the same hardware at once.
      \item \key{Data hazard:} an instruction depends on a value whose
        production or writeback has not completed.
      \item \key{Control hazard:} the next instruction address is uncertain,
        usually because of a branch or exception.
    \end{itemize}
  \end{block}
  \begin{itemize}
    \item The remedy must match the cause: add/duplicate a resource, move or
      forward data, or predict and recover from control flow.
  \end{itemize}
\end{frame}

\begin{frame}{Structural Hazard: One Port, Two Requests}
  \begin{exampleblock}{Conflict in cycle 4}
    In a five-stage pipeline, $I_1$ is in MEM while $I_2$ is in IF. If one
    single-ported memory serves both instruction fetches and data accesses,
    both stages request the same port in the same cycle.
  \end{exampleblock}
  \begin{itemize}
    \item \key{Option 1:} stall one instruction, inserting a bubble.
    \item \key{Option 2:} split instruction and data memories/caches, allowing
      the two accesses to proceed in parallel.
    \item \key{Option 3:} use multiported or time-multiplexed hardware, paying
      area, energy, or timing cost.
    \item A structural hazard is about \emph{resource demand}, not whether the
      instructions compute related values.
  \end{itemize}
\end{frame}

\begin{frame}{Data Hazards: The Running Stream}
  \begin{exampleblock}{RAW dependence}
    \texttt{LW R1,0(R2)}\\
    \texttt{ADD R3,R1,R4}\\
    \texttt{SUB R5,R3,R6}
  \end{exampleblock}
  \begin{itemize}
    \item \texttt{ADD} needs the value produced by \texttt{LW}; this is a
      \key{read after write (RAW)} dependence.
    \item \texttt{SUB} needs the value produced by \texttt{ADD}; another RAW
      dependence chains the stream.
    \item \key{WAR} (write after read) and \key{WAW} (write after write) can
      arise when instructions execute or complete out of order; in a simple
      in-order five-stage pipeline, RAW is the central hazard.
  \end{itemize}
\end{frame}

\begin{frame}{Timing Example 2: A Load-Use Stall}
  \scriptsize
  \begin{exampleblock}{Assumption}
    A load's data is available after MEM. The dependent ALU instruction needs
    it at the start of EX. One bubble is therefore required even with forwarding.
  \end{exampleblock}
  \begin{center}
    \begin{tabular}{@{}c|cccccccc@{}}
      \toprule
      & 1 & 2 & 3 & 4 & 5 & 6 & 7 & 8 \\
      \midrule
      \texttt{LW}  & IF & ID & EX & MEM & WB & & & \\
      \texttt{ADD} & & IF & ID & \textbf{stall} & EX & MEM & WB & \\
      \texttt{SUB} & & & IF & \textbf{stall} & ID & EX & MEM & WB \\
      \bottomrule
    \end{tabular}
  \end{center}
  \begin{itemize}
    \item The hazard unit freezes PC and IF/ID for one cycle and inserts a
      bubble into EX.
    \item The load's result can then be forwarded from the MEM/WB boundary to
      the dependent ADD's EX input.
  \end{itemize}
\end{frame}

\begin{frame}{Control Hazard: Which Instruction Is Next?}
  \begin{block}{The branch changes the instruction stream}
    \begin{itemize}
      \item The fetch unit may fetch the fall-through instruction before the
        branch condition is known.
      \item If the branch is taken, those fetched instructions are wrong-path
        work and must be prevented from changing architectural state.
      \item The penalty is the number of younger stages that must be flushed
        or redirected when the branch resolves.
    \end{itemize}
  \end{block}
  \begin{itemize}
    \item \key{Delay the decision:} simple, but inserts bubbles.
    \item \key{Predict and recover:} keep fetching; flush on a misprediction.
    \item \key{Move branch resolution earlier:} reduces penalty but lengthens
      an earlier stage and may add hardware to the critical path.
  \end{itemize}
\end{frame}

\begin{frame}{Socratic Check: Name the Hazard and the Cost}
  \begin{alertblock}{Three cases}
    \begin{enumerate}
      \item One memory port is requested by IF and MEM in the same cycle.
      \item An ADD reads a register before a preceding load produces it.
      \item A conditional branch has not yet resolved its target.
    \end{enumerate}
  \end{alertblock}
  \begin{itemize}
    \item (1) \key{Structural}: duplicate, split, or schedule the resource.
    \item (2) \key{Data / RAW}: forward when possible, otherwise stall.
    \item (3) \key{Control}: predict, delay, or flush wrong-path work.
    \item A complete answer states both the category and the mechanism that
      pays for avoiding it.
  \end{itemize}
\end{frame}

\transitionslide{Resolve the conflict at the cheapest point that preserves correctness.}

% =============================================================================
% ACT 4 --- RESOLUTION AND DESIGN TRADE-OFFS
% =============================================================================
\section{Hazard Resolution}

\begin{frame}{Resolution 1: Stall and Bubble}
  \begin{block}{The conservative mechanism}
    When the next instruction cannot safely advance, the hazard unit holds the
    blocked instructions and inserts a no-op bubble into the following stage.
  \end{block}
  \begin{itemize}
    \item The program's meaning is preserved because the dependent operation
      waits for the required resource or value.
    \item A stall can freeze the PC, freeze an inter-stage register, and clear
      the control fields entering the next stage.
    \item The cost is visible as extra cycles: if a fraction $s$ of instructions
      incur an average $b$ bubbles, then approximately
      $\mathrm{CPI}=1+s\times b$ before other effects.
  \end{itemize}
\end{frame}

\begin{frame}{Resolution 2: Forward the Value Before Writeback}
  \begin{exampleblock}{Why wait for the register file?}
    If an ALU result is already sitting at the EX/MEM boundary, the next ALU
    operation can consume it directly instead of waiting for WB to write and
    ID to read the register file.
  \end{exampleblock}
  \begin{itemize}
    \item Comparators detect whether a source register matches a destination
      register in a later pipeline stage.
    \item Multiplexers select the newest available value for each ALU input.
    \item Forwarding removes many ALU-to-ALU RAW stalls, but it cannot remove
      a load-use delay when memory data arrives too late.
  \end{itemize}
  \muted{Forwarding/data bypassing is part of the indexed pipeline treatment in
  Stallings 11e, \S16.4 and the RISC discussion in Ch.~17
  \cite{Stallings2015_computer_organization}.}
\end{frame}

\begin{frame}{Resolution 3: Schedule and Rename}
  \begin{itemize}
    \item \key{Compiler scheduling:} move an independent instruction into a
      delay slot or between a producer and consumer when the ISA permits it.
    \item \key{Register renaming:} give independent writes different physical
      destinations, removing false WAR and WAW name conflicts.
    \item \key{Scoreboarding/reservation stations:} track operand readiness
      and resource availability so independent instructions can proceed.
    \item These techniques exploit more ILP, but require metadata, selection
      logic, recovery rules, and stronger verification.
  \end{itemize}
  \begin{alertblock}{Boundary}
    No scheduling trick can remove a true RAW dependence: the consumer still
    needs the producer's value.
  \end{alertblock}
\end{frame}

\begin{frame}{Resolution 4: Predict, Then Recover}
  \begin{block}{Control-flow speculation}
    A branch predictor guesses direction and sometimes target so fetch continues
    without waiting for full branch resolution.
  \end{block}
  \begin{itemize}
    \item Correct prediction turns a control hazard into useful work.
    \item Wrong prediction requires a \key{flush}: invalidate younger
      instructions and restart from the correct target.
    \item Expected branch cost depends on branch frequency, misprediction rate,
      and recovery penalty:
      \[
        \Delta\mathrm{CPI}\approx f_{br}\,p_{miss}\,P_{flush}.
      \]
    \item Deeper pipelines can increase $P_{flush}$ even if their clock is faster.
  \end{itemize}
\end{frame}

\begin{frame}{Architecture Design: The Throughput Budget}
  \scriptsize
  \begin{center}
    \begin{tabular}{@{}p{2.8cm}p{3.2cm}p{3.2cm}p{3.0cm}@{}}
      \toprule
      \textbf{Choice} & \textbf{Potential gain} & \textbf{New cost} & \textbf{Question to ask} \\
      \midrule
      Deeper pipeline & Shorter stage delay; higher clock & More registers, bubbles, branch penalty & Are hazards predictable enough? \\
      Duplicate resources & Fewer structural conflicts & Area, power, coherence/consistency & Is the conflict frequent? \\
      Forwarding network & Fewer RAW stalls & Comparators, muxes, timing paths & Can the value arrive in time? \\
      Branch prediction & Higher control-flow throughput & Predictor state and flush recovery & What is the miss penalty? \\
      Out-of-order issue & More ILP exposed & Rename, queues, precise recovery & Is added complexity worth it? \\
      \bottomrule
    \end{tabular}
  \end{center}
  \vspace{0.1cm}
  \textbf{Design principle:} maximize useful completed work per unit time, not
  merely the clock rate or the number of stages.
\end{frame}

\begin{frame}{Running Stream: A Complete Hazard Diagnosis}
  \scriptsize
  \begin{exampleblock}{Stream}
    \texttt{LW R1,0(R2)}; \texttt{ADD R3,R1,R4}; \texttt{SUB R5,R3,R6};
    \texttt{BEQ R5,R0,L}; \texttt{OR R7,R8,R9}
  \end{exampleblock}
  \begin{itemize}
    \item \texttt{LW} $\to$ \texttt{ADD}: true RAW; forward plus one load-use
      stall in the five-stage timing model.
    \item \texttt{ADD} $\to$ \texttt{SUB}: true RAW; ALU forwarding can avoid
      an additional stall.
    \item \texttt{SUB} $\to$ \texttt{BEQ}: branch depends on a just-produced
      comparison value; forwarding or delayed branch resolution matters.
    \item \texttt{BEQ} $\to$ \texttt{OR}: control hazard; fetch may need a
      prediction and a flush if the branch is taken or mispredicted.
  \end{itemize}
  \muted{A timing diagram is a proof of the schedule: mark the producer stage,
  consumer stage, available bypass path, and every inserted bubble.}
\end{frame}

\begin{frame}{Socratic Check: Choose the Cheapest Correct Fix}
  \begin{alertblock}{Scenario}
    An ALU result is ready at the end of EX, its consumer enters EX next cycle,
    and the machine already has an ALU bypass path. What should the hazard unit
    do? What if the producer is a load instead?
  \end{alertblock}
  \begin{itemize}
    \item ALU producer: \good{forward the result}; no bubble is needed.
    \item Load producer: \key{stall until the data is available}, then forward;
      the memory access finishes too late for the immediately following EX.
    \item The correct resolution depends on \emph{when} the value becomes
      available, not only on the fact that a dependence exists.
  \end{itemize}
\end{frame}

\begin{frame}{Week 11 Synthesis: A Decision Procedure}
  \begin{enumerate}
    \item Identify the streams: instruction count and data count.
    \item For a pipeline conflict, name the blocked stage and requested resource
      or operand.
    \item Classify the hazard: structural, data, or control.
    \item Ask whether duplication, forwarding, scheduling, prediction, or a
      stall resolves it without changing program meaning.
    \item Count the bubbles or flushes and update the timing diagram.
    \item State the design trade-off: throughput, latency, area, power,
      complexity, and recovery cost.
  \end{enumerate}
  \begin{block}{Bridge to Week 12}
    Next week turns these mechanisms into worked instruction-pipelining
    problems, then compares RISC and CISC choices that make hazards easier or
    harder to manage.
  \end{block}
\end{frame}

\begin{frame}{References and Source Boundaries}
  \small
  \begin{itemize}
    \item Hamacher, Vranesic, Zaky, and Manjikian, \emph{Computer Organization
      and Embedded Systems}, 6th ed., Chs.~5 and 10. Chapter-level course
      anchor; current repository index does not page-verify these chapters.
    \item Stallings, \emph{Computer Organization and Architecture: Designing
      for Performance}, 11th ed. (repository key retained as Stallings2015),
      \S16.4, ``Instruction Pipelining,'' PDF pp.~667--685; \S20.1,
      ``Multiple Processor Organizations,'' PDF p.~847.
    \item Flynn's taxonomy is not attributed to Stallings here: its treatment
      in the indexed 11th edition is brief; the course anchor is Hamacher Ch.~5.
  \end{itemize}
  \vspace{0.2cm}
  \muted{Textbook page numbers are used only where the repository's index marks
  them as checkable. Chapter citations remain chapter citations.}
\end{frame}

\begin{frame}[allowframebreaks]{Bibliography}
  \bibliographystyle{plain}
  \bibliography{../../Bibliography_base}
\end{frame}

\end{document}